%=====================================================================
\chapter{The Data Science Landscape}\label{ch:landscape}
%=====================================================================
\locator{1}{Part I --- Foundations of Data}

\begin{outcomes}
By the end of this chapter you will be able to:
\begin{tight}
  \item \textbf{Define} data science and distinguish it from artificial intelligence, machine learning, deep learning, generative AI and agentic AI.
  \item \textbf{Explain} why these fields are nested rather than separate, using the containment diagram.
  \item \textbf{Classify} a problem into the level of analytics it requires --- descriptive, diagnostic, predictive or prescriptive.
  \item \textbf{Identify} the main roles on a data team and what each is accountable for.
  \item \textbf{Describe} one realistic data science application in each of the four running domains of this course.
\end{tight}
\end{outcomes}

\begin{prereq}
None. This is the entry point of the course. If the prerequisite self-check in the
front matter gave you trouble, keep Appendix C beside you.
\end{prereq}

\begin{keyterms}
data science $\bullet$ artificial intelligence $\bullet$ machine learning $\bullet$
deep learning $\bullet$ generative AI $\bullet$ agentic AI $\bullet$ data mining $\bullet$
descriptive, diagnostic, predictive and prescriptive analytics $\bullet$ data product
\end{keyterms}

\section{A Motivating Scenario}

A university notices that roughly one student in five fails the first-year
programming course. The department has years of records: enrolment forms, weekly
attendance, submission timestamps from the online judge, forum posts, quiz
scores. Nobody has ever looked at them together.

Three different questions could be asked of that pile of data, and they are not
the same question:

\begin{tight}
  \item \emph{How many students failed last year, and in which week did they stop submitting?} --- a question about what happened.
  \item \emph{Which students currently enrolled are likely to fail?} --- a question about what will happen.
  \item \emph{Which of these students should get a tutor, given that we can only fund fifteen?} --- a question about what to do.
\end{tight}

Data science is the discipline that answers all three, and that knows the
difference between them.

\section{What Data Science Is}

\begin{definitionbox}[Data Science]
Data science is the interdisciplinary practice of turning \textbf{data} into
\textbf{understanding} and \textbf{action}. It combines statistics, computing,
data engineering, machine learning, domain knowledge and communication to
extract insight, build predictive systems and support better decisions.
\end{definitionbox}

The definition is deliberately broad because the work is broad. Notice what it
contains besides mathematics: \emph{domain knowledge}, because a model of student
failure built by someone who has never taught will be wrong in ways nobody
catches; and \emph{communication}, because an insight nobody acts on has
accomplished nothing.

\begin{conceptbox}[The Three Things Data Science Needs at Once]
\emph{Mathematics and statistics} without programming leaves you unable to work
at scale. \emph{Programming} without statistics produces confident nonsense.
Both without \emph{domain knowledge} produce technically correct answers to the
wrong question. Most failures you will meet in industry are of the third kind.
\end{conceptbox}

\dsfig{%
\begin{tikzpicture}[scale=0.92]
  \begin{scope}[opacity=0.5, transparency group]
    \fill[secondary!50] (-1.5,0.9) circle (2.5cm);
    \fill[greenhl!50]   (1.5,0.9)  circle (2.5cm);
    \fill[goldhl!55]    (0,-1.3)   circle (2.5cm);
  \end{scope}
  \draw[secondary, line width=1.1pt] (-1.5,0.9) circle (2.5cm);
  \draw[greenhl,   line width=1.1pt] (1.5,0.9)  circle (2.5cm);
  \draw[goldhl,    line width=1.1pt] (0,-1.3)   circle (2.5cm);

  \node[font=\bfseries\small, secondary!50!black, align=center] at (-3.7,2.8) {Mathematics\\\& Statistics};
  \node[font=\bfseries\small, greenhl!45!black,  align=center] at ( 3.7,2.8) {Programming\\\& Engineering};
  \node[font=\bfseries\small, goldhl!45!black,   align=center] at ( 0,-4.2)  {Domain Knowledge};

  \node[font=\scriptsize, align=center, text=gray!45!black] at (0,1.62)     {Machine\\Learning};
  \node[font=\scriptsize, align=center, text=gray!45!black] at (-2.0,-0.72) {Traditional\\Research};
  \node[font=\scriptsize, align=center, text=gray!45!black] at (2.0,-0.72)  {Software\\Products};

  \node[font=\bfseries\footnotesize, align=center, text=primary,
        fill=white, rounded corners=3pt, inner sep=2.5pt] at (0,-0.1) {DATA\\SCIENCE};
\end{tikzpicture}}%
{Data science as the intersection of three capabilities. Each pair produces something
useful on its own; only the centre is data science. The most common professional failure
is drifting out of the gold circle --- solving a problem the domain did not have.}{ds-venn}

\section{The Nested Fields: AI, ML and Deep Learning}

Students routinely use ``AI'', ``machine learning'' and ``deep learning'' as
synonyms. They are not synonyms; they are nested. Getting this right in the first
week saves confusion for the rest of the course.

\dsfig{%
\begin{tikzpicture}[scale=0.9]
  \fill[black!8] (0.08,-0.08) circle (4.4cm);
  \fill[blue!6] (0,0) circle (4.4cm);
  \draw[thick, primary, line width=1.5pt] (0,0) circle (4.4cm);
  \node[font=\bfseries\small, primary] at (0,3.55) {\faIcon{robot}~Artificial Intelligence};
  \node[font=\scriptsize, primary!65] at (0,2.9)
       {rule-based systems $\bullet$ search $\bullet$ planning};

  \fill[orange!10] (0,-0.4) circle (2.85cm);
  \draw[thick, orange!70!black, line width=1.2pt] (0,-0.4) circle (2.85cm);
  \node[font=\bfseries\footnotesize, orange!70!black] at (0,1.85) {\faIcon{brain}~Machine Learning};
  \node[font=\scriptsize, orange!55!black] at (0,1.2)
       {regression $\bullet$ trees $\bullet$ clustering};

  \fill[green!12] (0,-0.85) circle (1.6cm);
  \draw[thick, greenhl!70!black, line width=1.2pt] (0,-0.85) circle (1.6cm);
  \node[font=\bfseries\scriptsize, greenhl!55!black, align=center] at (0,-0.45)
       {\faIcon{project-diagram}\\Deep Learning};
  \node[font=\tiny, greenhl!45!black] at (0,-1.5) {CNNs $\bullet$ Transformers};

  % the strictness of the containment, stated on the figure itself
  \node[font=\scriptsize\itshape, text=gray!55!black, align=center] at (0,-3.55)
       {every inner ring is a special case of the ring around it};
\end{tikzpicture}}%
{The nesting of AI, machine learning and deep learning. Every deep learning system is a
machine learning system; every machine learning system is an AI system; the reverse is
false in both cases. A hand-written rulebook that decides loan applications is AI but not
machine learning.}{ai-nesting}

\begin{tight}
  \item \term{Artificial Intelligence} is the outermost field: any technique that
        makes a machine behave in a way we would call intelligent. A chess program
        that searches moves exhaustively is AI and contains no learning at all.
  \item \term{Machine Learning} is the subset of AI in which behaviour is
        \emph{learned from data} rather than programmed. Nobody writes the rule
        ``if the email contains the word \emph{lottery}, mark it spam'' --- the
        algorithm discovers it.
  \item \term{Deep Learning} is the subset of ML that uses neural networks with
        many layers, which lets the system learn its own features instead of being
        handed them.
\end{tight}

\subsection{Where Data Science Sits}

Data science is not a fourth ring inside those circles --- it overlaps them.
Plenty of valuable data science involves no machine learning at all: a well-built
dashboard showing a hospital its live bed occupancy changes decisions every day
and contains not one model.

\dsfig{%
\begin{tikzpicture}[
  b/.style={rounded corners=3pt, draw=#1, line width=1pt, fill=#1!8, text=#1!50!black,
            align=center, font=\scriptsize\bfseries, minimum width=2.7cm, minimum height=0.9cm},
  arr/.style={-{Stealth[length=2.2mm]}, primary!50, line width=0.85pt}]

  \node[b=primary]   (ds)  at (0,0)      {Data Science};
  \node[b=secondary] (an)  at (-4.2,-1.8){Analytics \&\\Visualisation};
  \node[b=greenhl]   (ml)  at (-1.4,-1.8){Machine\\Learning};
  \node[b=goldhl]    (de)  at (1.4,-1.8) {Data\\Engineering};
  \node[b=purplehl]  (st)  at (4.2,-1.8) {Statistics \&\\Experiments};

  \foreach \n in {an,ml,de,st}{\draw[arr] (ds) -- (\n);}

  \node[b=greenhl, minimum width=2.3cm] (dl)  at (-1.4,-3.5) {Deep Learning};
  \node[b=greenhl, minimum width=2.3cm] (gen) at (-4.2,-3.5) {Generative AI};
  \node[b=greenhl, minimum width=2.3cm] (ag)  at (1.4,-3.5)  {Agentic AI};
  \draw[arr] (ml) -- (dl);
  \draw[arr] (dl) -- (gen);
  \draw[arr] (dl) -- (ag);
  \node[font=\scriptsize\itshape, text=gray!55!black, align=left, text width=3.5cm]
       at (4.6,-3.5) {Chapters 21--22 return to these once the neural network
       machinery of Chapter 17 is in place.};
\end{tikzpicture}}%
{Data science and its component practices. Analytics, engineering and statistics are
first-class members, not lesser cousins of machine learning --- most working days involve
more of them than of modelling.}{ds-components}

\section{The Newer Members: Generative and Agentic AI}

Two terms have entered everyday use recently and are worth pinning down now, even
though the course does not treat them properly until Part IV.

\begin{definitionbox}[Generative AI]
Systems that \textbf{create new content} --- text, images, code, audio --- resembling
their training data, rather than producing a label or a number. The question they
answer is \emph{``what would plausibly come next?''}
\end{definitionbox}

\begin{definitionbox}[Agentic AI]
Systems that pursue a \textbf{goal} over multiple steps: they plan, call tools,
observe results and adapt. The question they answer is \emph{``what should I do
next to get this done?''}
\end{definitionbox}

\begin{conceptbox}[Predictive, Generative, Agentic in One Sentence Each]
A \emph{predictive} model tells you a student is likely to fail. A
\emph{generative} model drafts the email you send them. An \emph{agentic} system
checks the timetable, finds a free tutor slot, drafts the email, books the
appointment and reports back. Each is strictly more capable --- and strictly
harder to keep under control.
\end{conceptbox}

\section{The Four Levels of Analytics}

The three questions in the opening scenario were not arbitrary. They correspond to
a standard ladder of analytical maturity, and one of the most useful things you can
do at the start of any project is work out which rung you are actually on.

\dsfig{%
\begin{tikzpicture}[
  step/.style={rounded corners=3pt, fill=#1, text=white, align=center,
               font=\bfseries\scriptsize, minimum width=2.7cm, minimum height=0.95cm},
  q/.style={font=\scriptsize\itshape, text=gray!55!black, align=center}]

  \draw[-{Stealth[length=3mm]}, primary!30, line width=7pt, opacity=0.4]
        (-0.4,-0.5) -- (11.4,3.05);

  \node[step=secondary] (d)  at (0.6,0)    {Descriptive\\\faIcon{chart-bar}};
  \node[step=tealhl]    (di) at (4.0,1.0)  {Diagnostic\\\faIcon{search}};
  \node[step=greenhl]   (p)  at (7.4,2.0)  {Predictive\\\faIcon{chart-line}};
  \node[step=purplehl]  (pr) at (10.8,3.0) {Prescriptive\\\faIcon{lightbulb}};

  \node[q] at (0.6,-1.1)  {What happened?};
  \node[q] at (4.0,-0.1)  {Why did it happen?};
  \node[q] at (7.4,0.9)   {What will happen?};
  \node[q] at (10.8,1.9)  {What should we do?};

  \node[font=\scriptsize, text=primary!60!black, rotate=17] at (5.4,3.15)
       {increasing value~~$\longrightarrow$~~increasing difficulty};
\end{tikzpicture}}%
{The four levels of analytics. Value rises as you climb, but so does the cost of being
wrong: a mistaken description misleads, a mistaken prescription causes harm.}{analytics-levels}

\rowcolors{2}{white}{lightbg}
\begin{longtable}{@{}L{2.4cm}L{4.0cm}L{8.2cm}@{}}
\toprule
\rowcolor{primary!15}
\textbf{Level} & \textbf{Typical technique} & \textbf{Example from the opening scenario} \\
\midrule
\endhead
Descriptive & Counts, averages, dashboards, visualisation & 21\% of students failed; submissions collapse in week 6 \\
Diagnostic & Segmentation, correlation, drill-down, root cause & Failures concentrate among students who missed the week 5 lab \\
Predictive & Classification, regression, forecasting & This student has an estimated 0.78 probability of failing \\
Prescriptive & Optimisation, simulation, decision rules & Assign the fifteen tutor slots so as to maximise expected passes \\
\bottomrule
\end{longtable}

\begin{pitfall}
\begin{tight}
  \item \textbf{Selling prescriptive, delivering descriptive.} The most common
        failure of student and industry projects alike. Be honest about which
        rung you reached.
  \item \textbf{Skipping straight to prediction.} If nobody has described the data
        yet, your model rests on a misunderstanding. Every hour spent on
        \chref{eda} saves three in \chref{evaluation}.
  \item \textbf{Assuming prescription follows from prediction.} Knowing who will
        fail does not tell you what helps. That needs the causal thinking of
        \chref{causality}.
\end{tight}
\end{pitfall}

\section{Who Does What: Roles on a Data Team}

\rowcolors{2}{white}{defbg}
\begin{longtable}{@{}L{3.2cm}L{5.4cm}L{5.8cm}@{}}
\toprule
\rowcolor{primary!15}
\textbf{Role} & \textbf{Accountable for} & \textbf{Chapters that matter most} \\
\midrule
\endhead
Data Analyst & Describing and diagnosing; dashboards and reports decision-makers actually use & 7, 8, 9, 31 \\
Data Scientist & Framing problems, building and validating models, quantifying uncertainty & 9--16, 28 \\
Data Engineer & Reliable pipelines; data arriving complete, on time and in the right shape & 3, 7, 23, 24 \\
ML Engineer & Getting models into production and keeping them correct there & 17--19, 27 \\
Security Analyst & Detecting and investigating threats in system and network data & 14, 16, 19, 26 \\
IoT / Edge Engineer & Sensor fleets, telemetry pipelines, on-device inference & 19, 24, 25 \\
Product / Project Manager & Choosing what is worth building, and delivering it & 2, 30, 31 \\
\bottomrule
\end{longtable}

\begin{conceptbox}[These Are Roles, Not People]
In a large organisation these are seven job titles. In a small one they are seven
hats worn by two people, and in your final-year project they will all be worn by
you. That is precisely why this handout covers all of them.
\end{conceptbox}

%---------------------------------------------------------------------
\begin{fourlenses}{The Levels of Analytics}
\lensLA{A university wants to reduce first-year attrition. \emph{Descriptive:}
dashboards of attendance and submission rates by cohort. \emph{Diagnostic:} which
combinations of missed activities precede withdrawal. \emph{Predictive:} a weekly
risk score per student. \emph{Prescriptive:} an allocation of limited tutoring
hours to the students it will help most. Only the first rung is reachable without
the rest of this book --- and only the first is safe to deploy without reading
\chref{ethics}.}

\lensPM{A software house delivers projects late about half the time.
\emph{Descriptive:} historical burn-down and estimate-versus-actual charts.
\emph{Diagnostic:} which task types are systematically underestimated.
\emph{Predictive:} a sprint-completion forecast published each Wednesday.
\emph{Prescriptive:} a recommended re-allocation of two developers. The predictive
rung is where project management stops being reporting and becomes data science.}

\lensIOT{A factory runs 400 motors instrumented for vibration and temperature.
\emph{Descriptive:} live telemetry dashboards. \emph{Diagnostic:} which sensor
pattern preceded last quarter's three failures. \emph{Predictive:} remaining useful
life per motor. \emph{Prescriptive:} a maintenance schedule minimising downtime
cost. Because data arrives continuously and in volume, even the descriptive rung
needs the engineering of \chref{dataeng}.}

\lensSEC{A bank monitors authentication logs across 12{,}000 accounts.
\emph{Descriptive:} failed-login counts by region and hour. \emph{Diagnostic:}
which accounts share an unusual user-agent string. \emph{Predictive:} a per-session
compromise likelihood. \emph{Prescriptive:} automatic step-up authentication for the
riskiest 0.1\% of sessions. Here an adversary is actively trying to look normal,
which is what makes the security lens harder than the other three.}
\end{fourlenses}

%---------------------------------------------------------------------
\begin{worked}[Classifying a Request]
A programme director asks: \emph{``Can you tell me whether the new lab format is
worth keeping?''} Which rung is really being requested?

\smallskip
\textbf{Step 1 --- restate it as something measurable.} ``Worth keeping'' must
become countable. Say: \emph{does the new lab format raise the end-of-semester
pass rate?}

\textbf{Step 2 --- identify the rung.} It is not descriptive (pass rates alone
would not settle it) and not predictive (nobody is asking about a future student).
It is \emph{diagnostic}, and specifically \emph{causal}: did the format cause the
change?

\textbf{Step 3 --- name what that requires.} A causal claim needs a comparison
group and a design that rules out rival explanations --- the machinery of
\chref{causality}. Comparing this year's pass rate with last year's confounds the
format with a different cohort, different staff and a different exam paper.

\textbf{Answer:} the request is a causal question disguised as a descriptive one.
Delivering a bar chart of pass rates would answer the wrong question convincingly,
which is worse than answering nothing.
\end{worked}

%---------------------------------------------------------------------
\begin{lab}[Meeting a Dataset]
The two most important facts about any new dataset are its \emph{shape} and its
\emph{types}. Every project in this course starts here.

\begin{lstlisting}[language=Python]
import pandas as pd

# A tiny extract of the kind of record a learning analytics project starts from.
df = pd.DataFrame({
    "student_id":     [101, 102, 103, 104, 105],
    "attendance_pct": [92, 45, 78, 30, 88],
    "quizzes_done":   [8, 3, 6, 2, 8],
    "forum_posts":    [12, 0, 4, 1, 20],
    "final_grade":    ["A", "F", "C", "F", "A"],
})

print(df.shape)        # (rows, columns) -> 5 observations, 5 features
print(df.dtypes)       # what type is each column?
print(df.describe())   # centre and spread of the numeric columns
\end{lstlisting}

\textbf{Then answer, without fitting anything:} which columns look related to
\texttt{final\_grade}? That informal glance is descriptive analytics, and it is
where every project should begin.
\end{lab}

%---------------------------------------------------------------------
\begin{checkpoint}
\begin{tightnum}
  \item A hospital deploys a hand-written rulebook: ``if temperature $>39^\circ$C
        and white cell count is high, flag for review.'' Is this AI? Is it machine
        learning? Justify both answers.
  \item Classify each request by analytics level: (a) ``How many tickets did we
        close last month?'' (b) ``Which server will run out of disk first?''
        (c) ``Why did response times spike on Tuesday?'' (d) ``What is the
        cheapest way to meet our uptime target?''
  \item A colleague says ``we're doing AI'' about a dashboard of monthly sales
        totals. What are they actually doing, and why does the distinction matter
        to whoever is paying for it?
\end{tightnum}
\end{checkpoint}

%---------------------------------------------------------------------
\begin{chaptersummary}
\begin{tight}
  \item Data science turns data into understanding and action; it needs
        mathematics, programming and domain knowledge \emph{simultaneously}.
  \item AI $\supset$ machine learning $\supset$ deep learning. The containment is
        strict in both directions of the argument: not all AI learns, and not all
        learning is deep.
  \item Generative AI creates content; agentic AI pursues goals over multiple
        steps. Both build on the deep learning of Part IV.
  \item Analytics has four rungs --- descriptive, diagnostic, predictive,
        prescriptive. Value and risk both increase as you climb.
  \item Much valuable data science contains no model at all. Do not mistake
        modelling for the whole discipline.
\end{tight}
\end{chaptersummary}

\begin{reviewq}
  \item \textbf{(MCQ)} Which statement is true?
        (a)~All machine learning is deep learning.
        (b)~All deep learning is machine learning.
        (c)~AI and machine learning are the same field.
        (d)~Data science is a subset of deep learning.
  \item \textbf{(MCQ)} A model estimates the probability that a motor fails within
        30 days. This is:
        (a)~descriptive (b)~diagnostic (c)~predictive (d)~prescriptive.
  \item \textbf{(Short)} Give one example of an AI system that contains no machine
        learning, and explain what makes it AI nonetheless.
  \item \textbf{(Short)} Explain, in two sentences, the difference between a
        generative and an agentic system.
  \item \textbf{(Short)} Why is domain knowledge treated in this chapter as a
        capability equal to statistics and programming rather than as background?
  \item \textbf{(Applied)} A retailer asks you to ``use AI to increase sales.''
        Rewrite this as one question at each of the four analytics levels, and
        state which you would deliver first and why.
  \item \textbf{(Applied)} For each of the four running domains of this course,
        name one descriptive question and one predictive question that a
        first-year BSIT student could realistically attempt.
\end{reviewq}
